\begin{proofbox}
\textbf{方法一：平方差配合距离差（优先）。}
% P=(15/4,3), H=(15/4,0); uniform units, a=3,b=4,c=5.
\begin{center}\begin{tikzpicture}[font=\small,>=stealth]
\draw[black!30,rounded corners=2pt] (0,0) rectangle (14.4,8.2);
\node[anchor=west,font=\small\bfseries] at (.5,7.72) {两条焦半径共用一条高：平方相减时，高的平方消掉};
\begin{scope}[shift={(6.9,3.5)},scale=.78]
\draw[->,black!40] (-6,0)--(6.5,0) node[right] {$x$};
\draw[blue,very thick,domain=0:4.1,samples=101,variable=\t]
 plot ({3*sqrt(1+\t*\t/16)},\t);
\draw[orange!85!black,thick] (-5,0)--(3.75,3)
 node[pos=.43,above=6pt] {$r_1=9.25$};
\draw[orange!85!black,thick] (5,0)--(3.75,3)
 node[midway,right=9pt] {$r_2=3.25$};
\draw[violet,thick,dashed] (3.75,3)--(3.75,0)
 ;
\node[violet] at (1.1,.8) {$PH=3$};
\draw[violet] (3.47,0)--(3.47,.28)--(3.75,.28);
\fill (-5,0) circle (2pt) node[above left=4pt] {$F_1(-5,0)$};
\fill (5,0) circle (2pt) node[below right=4pt] {$F_2(5,0)$};
\fill (3.75,0) circle (2pt) node[below left=5pt] {$H$};
\fill (3.75,3) circle (2pt) node[above left=5pt] {$P(3.75,3)$};
\draw[<->,black!60] (-5,-.85)--(3.75,-.85)
 node[midway,below=3pt] {$F_1H=8.75$};
\draw[<->,black!60] (3.75,-1.75)--(5,-1.75)
 node[midway,below=3pt] {$HF_2=1.25$};
\end{scope}
\node[anchor=west] at (.7,1.05) {$r_1^2-r_2^2=(8.75^2+3^2)-(1.25^2+3^2)=75=4\cdot5\cdot3.75$};
\node[anchor=west] at (.7,.5) {推广：$(x+c)^2+y^2-\bigl((x-c)^2+y^2\bigr)=4cx$；不必先解出 $y$};
\end{tikzpicture}\end{center}

图中 $P=(15/4,3)$、$H=(15/4,0)$，两条焦半径是两个直角三角形的斜边。
它们共用高 $PH$，所以勾股式相减时同一个高度平方消去。
图中虽有 $a<x<c$，下面的 $(x\pm c)^2$ 已自动包含水平距离的绝对值，
故推导对 $x\ge c$ 及下半支也适用。
\par\medskip\noindent
设 $F_1(-c,0),F_2(c,0)$，$r_1=PF_1,r_2=PF_2$。
右支取 $r_1-r_2=2a$，距离公式给出
\[
(r_1-r_2)(r_1+r_2)=4cx.
\]
所以 $r_1+r_2=2cx/a$，进而
\[
r_1=\frac{cx}{a}+a,\qquad r_2=\frac{cx}{a}-a.
\]
把第二式代入 $r_2^2=(x-c)^2+y^2$，得
\[
\frac{c^2x^2}{a^2}+a^2=x^2+c^2+y^2,
\qquad (c^2-a^2)x^2-a^2y^2=a^2(c^2-a^2).
\]
除以 $a^2b^2>0$ 得标准方程。左支关于 $y$ 轴对称，也得到相同方程。

\textbf{反向必须分支取正根。}
由标准方程有 $|x|\ge a$，且
\[
r_1^2=\left(\frac{cx}{a}+a\right)^2,\qquad
r_2^2=\left(\frac{cx}{a}-a\right)^2.
\]
若 $x\ge a$，两个括号均正，因为 $cx/a-a\ge c-a>0$，故 $r_1-r_2=2a$。
若 $x\le-a$，两个括号均负，取距离的正平方根后得到 $r_1-r_2=-2a$。
所以 $|r_1-r_2|=2a$，等价性成立。
\end{proofbox}
\clearpage
\begin{proofbox}
\textbf{方法二：直接移项平方。}
由 $r_1=r_2\pm2a$，平方并用 $r_1^2-r_2^2=4cx$，得
\[
cx-a^2=\pm ar_2.
\]
再平方为 $(cx-a^2)^2=a^2((x-c)^2+y^2)$，整理即得标准方程。
这种写法较熟悉，但反向仍需方法一的分支验证，平方本身不能保存距离差的符号。

\textbf{几何角度：三种边界。}
由 $PF_1\le PF_2+F_1F_2$ 及交换焦点后的不等式，
\[
|PF_1-PF_2|\le F_1F_2.
\]
差为零时，两距离相等，点在垂直平分线上；差等于焦距时，只能共线且位于焦点线段外，含两个焦点。
\begin{center}\begin{tikzpicture}[font=\small,>=stealth]
\draw[black!30,rounded corners=2pt] (0,0) rectangle (14.4,8.2);
\node[anchor=west,font=\small\bfseries] at (.5,7.72) {固定焦距 $10$：差为零与差等于焦距，都不是双曲线};
\node at (3.6,6.65) {$|PF_1-PF_2|=0$};
\node at (10.8,6.65) {$|PF_1-PF_2|=10$};
\begin{scope}[shift={(3.6,3.8)},scale=.32]
\draw[black!40] (-8,0)--(8,0);
\draw[blue,very thick] (0,-7)--(0,7);
\draw[orange!85!black] (-5,0)--(0,4)--(5,0);
\fill (-5,0) circle (2pt) node[below] {$F_1$};
\fill (5,0) circle (2pt) node[below] {$F_2$};
\fill (0,4) circle (2pt) node[above right] {$P$};
\end{scope}
\begin{scope}[shift={(10.8,3.8)},scale=.32]
\draw[black!40] (-8,0)--(8,0);
\draw[blue,very thick,->] (-5,0)--(-8,0);
\draw[blue,very thick,->] (5,0)--(8,0);
\fill (-5,0) circle (2pt) node[below] {$F_1$};
\fill (5,0) circle (2pt) node[below] {$F_2$};
\end{scope}
\node at (3.6,1.3) {垂直平分线};\node at (10.8,1.3) {两条向外射线，含焦点};

\node[anchor=west] at (.7,.55) {同尺度；任意点都有 $|PF_1-PF_2|\le10$，超过 $10$ 时轨迹为空集};
\end{tikzpicture}\end{center}

差大于焦距时为空集。正常的双曲线须排除以上边界，即 $0<2a<2c$。
\end{proofbox}
\clearpage

% ===== C037 扩充区开始：原有方法一、方法二及边界证明全部保留 =====
\begin{proofbox}
\textbf{补充方法三：考场“最少步骤”版——差平方，一消一留。}
设 $F_1(-c,0)$、$F_2(c,0)$，动点 $P(x,y)$。
若在右支，记 $r_1=PF_1,r_2=PF_2$，且 $r_1-r_2=2a$。
\textbf{第一步：} 展开两条距离的平方，只保留相减的结果（$y^2$ 自动消掉）：
\[
\underbrace{\bigl((x+c)^2+y^2\bigr)}_{r_1^2}
-\underbrace{\bigl((x-c)^2+y^2\bigr)}_{r_2^2}=4cx.
\]
\textbf{第二步：} 利用\textbf{平方差 $=$ 差 $\times$ 和}，
\[
(r_1-r_2)(r_1+r_2)=4cx
\quad\Longrightarrow\quad r_1+r_2=\frac{2cx}{a}.
\]
一加一减解出
\[
\boxed{r_1=\frac{cx}{a}+a,\qquad r_2=\frac{cx}{a}-a}\quad(x\ge a).
\]
\textbf{第三步：} 用第二个距离的勾股式还原 $y^2$：
\[
\begin{aligned}
\left(\frac{cx}{a}-a\right)^2&=(x-c)^2+y^2,\\
\frac{c^2x^2}{a^2}-2cx+a^2
 &=x^2-2cx+c^2+y^2,\\
y^2&=\frac{c^2-a^2}{a^2}x^2-(c^2-a^2)
   =\frac{b^2}{a^2}(x^2-a^2).
\end{aligned}
\]
由 $b^2>0$ 立即得到 $x^2/a^2-y^2/b^2=1$。
\textbf{速记顺序：两距离平方相减 $\to$ 平方差 $\to$ 解出一条距离 $\to$ 勾股还原。}
这和方法一是同一原理，只是将容易漏写的\emph{中间桥梁}独立展开。
\end{proofbox}
\clearpage
\begin{proofbox}
\textbf{补充方法四：先用数轴找顶点，再证明分支为何从顶点开始。}
沿焦点所在的 $x$ 轴，取 $-c\le x\le c$，有
\[
PF_1=x+c,\qquad PF_2=c-x,
\quad |PF_1-PF_2|=2|x|.
\]
令 $2|x|=2a$，便得两个顶点 $A_1(-a,0),A_2(a,0)$。
在右支，方法一已算出
\[
r_1+r_2=\frac{2cx}{a}.
\]
而三角形两边之和至少为焦距 $2c$，故
\[
\frac{2cx}{a}\ge 2c
\quad\Longrightarrow\quad x\ge a.
\]
这说明右支没有比右顶点更靠左的点；
左支镜像可得 $x\le-a$。
注意：此处等号对应 $P$ 位于焦点连线上，
所以\textbf{顶点也是焦点三角形允许退化的位置}。

\textbf{补充方法五：从第一定义顺手看见“第二定义”（选学）。}
方法一配合左右支的距离正号，可统一写为
\[
PF_2=\left|\frac{cx}{a}-a\right|
=\frac ca\left|x-\frac{a^2}{c}\right|.
\]
点 $P(x,y)$ 到直线 $x=a^2/c$ 的距离正是 $|x-a^2/c|$。
因此同一条双曲线也满足
\[
\boxed{\frac{PF_2}{d(P,\,x=a^2/c)}=\frac ca>1},
\]
其中 $d(P,\,x=a^2/c)$ 表示点到直线的距离。
这把熟悉的“点到直线距离”与双曲线联系起来；
初次学习可只掌握前四种角度，不必增加记忆负担。

\textbf{核验提醒：} 凡是平方或去掉绝对值后得到的方程，
必须再确认距离为非负，并验证左右支的符号；原文的反向验证应保留。
\end{proofbox}
\clearpage
